% Demonstration content only; replace it with sourced assignment content.
\documentclass[a4paper,14pt,times]{G7-32}
\usepackage{guap-report}
\begin{document}
\begin{GuapTitle}
\centering
{[Наименование организации из предоставленного титульного листа]}\par
{[Кафедра]}\par
\vfill
ОТЧЁТ\par
{[Вид работы, дисциплина и тема]}\par
\vfill
{[ФИО студента и группа]}\par
{[ФИО преподавателя]}\par
\vfill
{[Город и год]}\par
\end{GuapTitle}
\frontmatter
\tableofcontents
\chapter*{Введение}
\addcontentsline{toc}{chapter}{Введение}
Это демонстрация вёрстки, а не выполненная учебная работа. Содержание,
реквизиты и расчёты необходимо взять из задания и исходных материалов.
\mainmatter
\chapter{Основная часть}
\section{Исходные данные}
Демонстрационные значения приведены в таблице~\ref{tab:demo}.
Они не являются измерениями или результатами лабораторной работы.
\begin{table}[htbp]
\caption{Демонстрационные данные}\label{tab:demo}
\centering
\begin{tabular}{|p{90mm}|p{40mm}|}
\hline
Показатель & Значение \\
\hline
Демонстрационный показатель & 1 \\
\hline
\end{tabular}
\end{table}
\section{Расчёт}
Демонстрационная формула:
\begin{equation}
y = x + 1.
\end{equation}
Рисунки вставляйте через \texttt{includegraphics}, подпись размещайте снизу,
а ссылку на рисунок приводите в тексте перед изображением.
\backmatter
\chapter*{Заключение}
\addcontentsline{toc}{chapter}{Заключение}
Демонстрация содержит общую структуру. Состав разделов определяется заданием,
а не перечнем глав чужого отчёта.
\begin{thebibliography}{9}
\bibitem{gost} ГОСТ 7.32–2017. Система стандартов по информации,
библиотечному и издательскому делу. Отчёт о научно-исследовательской работе.
Структура и правила оформления. — Текст : непосредственный.
\end{thebibliography}
\end{document}
